\section{Algoritmo e implementación}\label{sec:algoritmo}

El modelo se resuelve en cinco etapas secuenciales, cada una consumiendo la salida de la anterior.

\textbf{(i) Datos del caso.} Los parámetros del caso asignado (presupuestos, límites del Cuadro~4 del caso, nivel de
confianza, aversión al riesgo, tamaño y semilla del remuestreo) y las posiciones y flujos contractuales
limpios se consolidan en un único documento de entrada, para que ningún límite, costo o posición quede fijado
fuera de esos datos.

\textbf{(ii) Valorización.} Cada instrumento se valora a $t_0$ descontando sus flujos contractuales con la
curva cero de su moneda más un spread calibrado para que ese valor coincida con el valor de mercado
observado, con una tolerancia de $10^{-6}$. La misma función de valorización, con
el mismo spread fijo, revaloriza cada instrumento bajo cualquier escenario de mercado.

\textbf{(iii) Escenarios y riesgo.} Se generan las 504 trayectorias descritas en la Sección~\ref{sec:riesgo} (500
por remuestreo más 4 de estrés) y se revaloriza cada uno de los doce instrumentos bajo cada una, para obtener
la matriz de resultados $r_{k,s}$ que alimenta la optimización. Calcular esa matriz completa toma cerca de un
minuto y se reutiliza en todas las variantes que no cambian los escenarios ni la valorización.

\textbf{(iv) Optimización.} Con $r_{k,s}$ fija, el problema lineal de la ecuación~\eqref{eq:objetivo} se
plantea con \textcite{diamond2016} y se resuelve con el solver de programación lineal de código abierto
HiGHS. La instancia base (12 instrumentos, 504 escenarios y unas 40 restricciones de peso y presupuesto)
converge en 167 iteraciones de simplex y 11 milisegundos. Cuando un grupo de peso queda vacío por la
composición de los instrumentos y su límite inferior es positivo, o cuando los límites de caja o patrimonio
resultan incompatibles entre sí, el problema se declara sin solución en vez de forzar un resultado.

\textbf{(v) Resultados.} La solución óptima, el barrido de la aversión al riesgo $\lambda$, un punto de
comparación de mínima varianza con el mismo resultado neto esperado, el diagnóstico de la posición inicial y
la variante que exige los límites de vencimiento también en el horizonte (Sección~\ref{sec:resultados}) se calculan
en la misma corrida y quedan disponibles como un único registro de resultados: cada cifra, tabla y figura de
este informe se lee directamente de ese registro, sin transcripción manual.

Todo el proceso es reproducible con una semilla fija para el remuestreo. La elección de revalorizar cada
escenario en vez de propagar la matriz de covarianzas de los 24 factores de riesgo (Sección~\ref{sec:datos}) evita
tener que reducir su dimensión: esa matriz es singular y una aproximación lineal exigiría antes una
descomposición de componentes principales o un filtro equivalente, con su propio error de aproximación.
